\chapter{Geometric Optics}

Every chapter so far has taken the wave nature of light seriously: interference, diffraction, polarization. Geometric optics is the opposite limit --- the $\lambda \to 0$ limit of the same theory --- in which light may be treated as travelling along \emph{rays}, the curves orthogonal to the wavefronts. Nothing new is postulated here. Every apparent peculiarity below (virtual images, signed distances, reversibility) is inherited from the wave theory already developed.

% ============================================================
\section{Domain of Validity, Terminology, and Sign Conventions}
% ============================================================

The single-slit result of Section~\ref{sec:single-slit} fixes exactly when ray language is safe. Light passing an aperture of size $a$ spreads by an angle $\sim\lambda/a$; only when every relevant dimension --- aperture, lens, mirror, object --- greatly exceeds $\lambda$ is that spreading negligible and the beam's edges sharp enough to call a ray. A lens $5\,$cm across passing visible light has $\lambda/a \sim 10^{-5}$, utterly negligible; the same light through a $10\,\mu$m pinhole spreads by several degrees, and ray optics fails outright.

The vocabulary is worth fixing before the derivations:
\begin{itemize}
    \item \textbf{Conjugate points}: an object point and its image point. By reversibility (below), the two labels may be exchanged.
    \item \textbf{Stigmatic imaging}: all rays leaving one object point pass exactly through one image point. Real surfaces achieve this only approximately, and only for particular conjugate pairs.
    \item \textbf{Vertex} $V$: where a refracting or reflecting surface meets the optical axis. All distances are measured from it.
    \item \textbf{Paraxial} (Gaussian, first-order) \textbf{optics}: the approximation keeping only the leading order in ray height and angle. Departures from it are \textbf{aberrations}.
    \item \textbf{Real image}: rays physically converge there, so it can be caught on a screen. \textbf{Virtual image}: rays only appear to diverge from it.
\end{itemize}

Every distance appearing below is signed, and one rule fixes them all: \textbf{positive means on the side where the light actually is}. Worth fixing now, since every equation from the next section onward depends on it.

\begin{center}
\begin{tabular}{lll}
\toprule
Quantity & Positive & Negative \\
\midrule
$s_o$ & object on incoming side & virtual object \\
$s_i$ & image on outgoing side (real) & image on incoming side (virtual) \\
$f$ (lens) & converging & diverging \\
$f$ (mirror) & concave & convex \\
$R$ & centre of curvature on outgoing side & otherwise \\
\bottomrule
\end{tabular}
\end{center}

A reliable anchor when memory fails: a biconvex lens in air must converge, so $f > 0$, which forces $R_1 > 0$ and $R_2 < 0$.

% ============================================================
\section{Foundational Principles}
% ============================================================

Section~\ref{sec:fermat} already stated Fermat's principle in its minimum-transit-time form and used it to derive Snell's law. The form used throughout this chapter is the equivalent statement in terms of \emph{optical path length},
\begin{formula}[Fermat's principle]
\begin{equation}
    \mathrm{OPL} = \int_A^B n(\ell)\,\dd\ell \quad \text{is stationary with respect to variations of the path.}
\end{equation}
\end{formula}
Since $n = c/v$, dividing the OPL by $c$ gives exactly the transit time, so this is the same principle. Note ``stationary,'' not ``minimum'': saddle points and maxima also correspond to physical rays.

The law of reflection $\theta_i = \theta_r$ and Snell's law $n_1\sin\theta_1 = n_2\sin\theta_2$ were both derived in Section~\ref{sec:refraction}, once from wavefront coherence and again from Fermat, and are simply used from here on. It is worth noticing that Snell's law can be written
\begin{equation}
    \frac{\sin\theta_1}{v_1} = \frac{\sin\theta_2}{v_2},
\end{equation}
in which $c$ has cancelled entirely: only the ratio $n_2/n_1 = v_1/v_2$ has physical content, and the law in this form governs \emph{any} wave --- sound, seismic --- not just light.

\begin{fact}[Reversibility]
Any ray traced backwards retraces its forward route exactly.
\end{fact}

This follows directly from Fermat: $\int n\,\dd\ell$ integrates a scalar against unsigned arc length, so it is unchanged under $A \leftrightarrow B$. Reversing each path is a bijection between the candidate paths of the two problems under which the functional takes equal values, so the stationary paths coincide. This is why the object and image labels of a conjugate pair may be swapped at will.

Total internal reflection was derived in Section~\ref{sec:tir}: for $n_1 > n_2$ no real refraction angle exists beyond $\sin\theta_c = n_2/n_1$. The wave picture adds what the ray picture cannot see: past $\theta_c$ the transmitted wavevector component normal to the interface becomes imaginary, so the field decays evanescently into the second medium, carries no time-averaged energy across it, and the reflection coefficient has $|r| = 1$ --- reflection is total in the strict sense, not merely strong.

% ============================================================
\section{Refraction at a Single Spherical Surface}
% ============================================================

Everything else in this chapter --- lenses, mirrors, compound instruments --- is built from one calculation: refraction at a single spherical surface, done exactly from Fermat's principle, with Snell's law never assumed.

The geometry is set up in Figure~\ref{fig:sphere-refract}. A surface of radius $R$ has centre $C$ and vertex $V$; an object $S$ sits a distance $s_o$ to the left in medium $n_1$, and we ask where the rays reconverge in medium $n_2$. A ray leaves $S$, meets the surface at $A$, and continues to $P$ at distance $s_i$. The single parameter locating $A$ is the angle $\varphi = \angle ACV$.

\begin{figure}[H]
\centering
\begin{tikzpicture}[>=Stealth, scale=0.95]
  \def\R{3}
  \def\phideg{38}
  \coordinate (V) at (0,0);
  \coordinate (C) at (\R,0);
  \coordinate (S) at (-5,0);
  \coordinate (P) at (7.6,0);
  \coordinate (A) at ({\R-\R*cos(\phideg)},{\R*sin(\phideg)});
  % media shading
  \fill[LightOrange,opacity=0.45] (-6,-2.8) -- (0,-2.8) --
      plot[domain=-55:55,samples=60] ({\R-\R*cos(\x)},{\R*sin(\x)}) -- (-6,2.8) -- cycle;
  \fill[LightGray,opacity=0.5] (8.6,-2.8) -- (0,-2.8) --
      plot[domain=-55:55,samples=60] ({\R-\R*cos(\x)},{\R*sin(\x)}) -- (8.6,2.8) -- cycle;
  % surface
  \draw[very thick] plot[domain=-55:55,samples=80] ({\R-\R*cos(\x)},{\R*sin(\x)});
  % axis
  \draw[dashed] (-6,0) -- (8.6,0);
  % rays
  \draw[very thick,blue!70!black,->] (S) -- ($(S)!0.62!(A)$);
  \draw[very thick,blue!70!black] (S) -- (A);
  \draw[very thick,red!75!black,->] (A) -- ($(A)!0.6!(P)$);
  \draw[very thick,red!75!black] (A) -- (P);
  % radius
  \draw[thick,dashed,gray] (C) -- (A);
  % angle phi at C
  \draw[thick] ($(C)+(180:0.85)$) arc (180:{180-\phideg}:0.85);
  \node at ($(C)+({180-\phideg/2}:1.15)$) {\small$\varphi$};
  % points
  \foreach \p/\n/\pos in {S/$S$/below,V/$V$/below left,C/$C$/below,P/$P$/below,A/$A$/above right}{
      \filldraw (\p) circle (1.6pt); \node[\pos] at (\p) {\small\n};
  }
  % length labels
  \node[blue!70!black] at (-2.9,1.55) {\small$\ell_o$};
  \node[red!75!black] at (4.2,1.5) {\small$\ell_i$};
  \node[gray] at (2.55,1.45) {\small$R$};
  % distances
  \draw[<->] (-5,-2.35) -- (0,-2.35) node[midway,fill=LightOrange,inner sep=1pt] {\small$s_o$};
  \draw[<->] (0,-2.75) -- (7.6,-2.75) node[midway,fill=LightGray,inner sep=1pt] {\small$s_i$};
  \node at (-4.6,2.3) {\small$n_1$};
  \node at (7.6,2.3) {\small$n_2$};
\end{tikzpicture}
\caption{Refraction at a single spherical surface. The ray $S \to A \to P$ is located by the single angle $\varphi = \angle ACV$; Fermat's principle demands that the optical path $n_1\ell_o + n_2\ell_i$ be stationary as $A$ slides along the surface.}
\label{fig:sphere-refract}
\end{figure}

Applying the law of cosines to triangles $SAC$ and $ACP$,
\begin{align}
    \ell_o &= \sqrt{R^2 + (s_o+R)^2 - 2R(s_o+R)\cos\varphi}, \\
    \ell_i &= \sqrt{R^2 + (s_i-R)^2 + 2R(s_i-R)\cos\varphi},
\end{align}
the sign difference arising because the angle included at $C$ is $\varphi$ in the first triangle and $\pi - \varphi$ in the second. Fermat requires $\mathrm{OPL} = n_1\ell_o + n_2\ell_i$ to be stationary as the ray's landing point slides along the surface:
\begin{equation}
    \dv{\varphi}\left(n_1\ell_o + n_2\ell_i\right) = 0
    \;\Longrightarrow\;
    \frac{n_1 R(s_o+R)\sin\varphi}{\ell_o} - \frac{n_2 R(s_i-R)\sin\varphi}{\ell_i} = 0.
\end{equation}
Dividing through by $R\sin\varphi$ and rearranging gives an \emph{exact} relation, with no approximation anywhere:
\begin{formula}[Exact single-surface relation]
\begin{equation}
    \frac{n_1}{\ell_o} + \frac{n_2}{\ell_i}
    = \frac{1}{R}\left(\frac{n_2 s_i}{\ell_i} - \frac{n_1 s_o}{\ell_o}\right).
\end{equation}
\end{formula}

For rays close to the axis the surface sagitta is negligible, so $\ell_o \to s_o$ and $\ell_i \to s_i$, and the relation collapses to the workhorse of paraxial optics:
\begin{formula}[Paraxial single-surface equation]
\begin{equation}
    \frac{n_1}{s_o} + \frac{n_2}{s_i} = \frac{n_2 - n_1}{R}.
\end{equation}
\end{formula}

The ray height has dropped out entirely, and that is the whole reason an image forms at all: \emph{every} paraxial ray leaving $S$ arrives at the same $P$, independent of where it struck the surface. Retaining the next order in $\varphi$ restores a height dependence, and that residual is precisely spherical aberration.

Setting $s_i \to \infty$ and $s_o \to \infty$ in turn gives the two focal lengths
\begin{equation}
    f_o = \frac{n_1 R}{n_2 - n_1}, \qquad
    f_i = \frac{n_2 R}{n_2 - n_1}, \qquad
    \frac{f_i}{f_o} = \frac{n_2}{n_1},
\end{equation}
which differ precisely because the media differ; they coincide only when the surrounding medium is the same on both sides, as for a lens in air.

\begin{example}[Two limiting checks]
Setting $R \to \infty$ flattens the surface into a plane interface and leaves $s_i = -(n_2/n_1)s_o$: a virtual image at reduced apparent depth. For water viewed from air, $n_1 = 1.33$ and $n_2 = 1$, so the bottom of a pool appears at about $0.75$ of its true depth. Setting $n_1 = n_2$ instead gives $s_i = -s_o$, meaning the ``image'' sits exactly on the object --- correct, since with no index contrast there is no surface at all.
\end{example}

% ============================================================
\section{Thin Lenses}
% ============================================================

A lens is two spherical surfaces in series, and nothing beyond the paraxial single-surface equation is needed to handle it. The construction is shown in Figure~\ref{fig:lensmaker}: a lens of index $n_l$ and thickness $d$, with vertices $V_1$ and $V_2$ and radii $R_1$ and $R_2$, sits in a surrounding medium $n_m$.

\begin{figure}[H]
\centering
\begin{tikzpicture}[>=Stealth, scale=1.05]
  \def\Ra{3.5}
  % lens body: two arcs, centres at (3.05,0) and (-3.05,0)
  \fill[blue!8]
      plot[domain=153.71:206.29,samples=60] ({3.05+\Ra*cos(\x)},{\Ra*sin(\x)}) --
      plot[domain=-26.29:26.29,samples=60] ({-3.05+\Ra*cos(\x)},{\Ra*sin(\x)}) -- cycle;
  \draw[very thick] plot[domain=153.71:206.29,samples=60] ({3.05+\Ra*cos(\x)},{\Ra*sin(\x)});
  \draw[very thick] plot[domain=-26.29:26.29,samples=60] ({-3.05+\Ra*cos(\x)},{\Ra*sin(\x)});
  % axis
  \draw[dashed] (-8.3,0) -- (7.1,0);
  % key points
  \coordinate (S)  at (-5,0);
  \coordinate (Pp) at (-7.5,0);
  \coordinate (P)  at (6,0);
  \coordinate (A)  at (-0.289,1.05);
  \coordinate (B)  at (0.262,1.13);
  % ray: S -> A (in n_m), A -> B (in n_l), B -> P (in n_m)
  \draw[very thick,blue!70!black,->] (S) -- ($(S)!0.75!(A)$);
  \draw[very thick,blue!70!black] (S) -- (A);
  \draw[very thick,Green] (A) -- (B);
  \draw[very thick,red!75!black,->] (B) -- ($(B)!0.72!(P)$);
  \draw[very thick,red!75!black] (B) -- (P);
  % backward extension of the in-lens ray to the virtual intermediate image
  \draw[thick,Green,dashed] (A) -- (Pp);
  % labelled points
  \foreach \p/\n/\pos in {S/$S$/below,Pp/$P'$/below,P/$P$/below}{
      \filldraw (\p) circle (2pt); \node[\pos] at (\p) {\small\n};
  }
  \filldraw (-0.45,0) circle (1.6pt); \node[above left] at (-0.47,0.02) {\small$V_1$};
  \filldraw (0.45,0) circle (1.6pt);  \node[above right] at (0.47,0.02) {\small$V_2$};
  \filldraw (3.05,0) circle (1.6pt);  \node[below] at (3.05,-0.08) {\small$C_1$};
  \filldraw (-3.05,0) circle (1.6pt); \node[below] at (-3.05,-0.08) {\small$C_2$};
  % media
  \node at (-4.3,1.9) {\small$n_m$};
  \node at (3.9,1.9) {\small$n_m$};
  \node at (0,-0.85) {\scriptsize$n_l$};
  % distance brackets
  \draw[<->] (-0.45,-2.05) -- (0.45,-2.05);
  \node[below] at (0,-2.07) {\small$d$};
  \draw[<->] (-5,-2.65) -- (-0.45,-2.65) node[midway,fill=white,inner sep=1pt] {\small$s_{o1}$};
  \draw[<->] (0.45,-2.65) -- (6,-2.65) node[midway,fill=white,inner sep=1pt] {\small$s_{i2}$};
  \draw[<->] (-7.5,-3.2) -- (-0.45,-3.2) node[midway,fill=white,inner sep=1pt] {\small$|s_{i1}|$};
  \draw[<->] (-7.5,-3.75) -- (0.45,-3.75) node[midway,fill=white,inner sep=1pt] {\small$s_{o2}$};
  % thin drop lines
  \foreach \x in {-7.5,-5,-0.45,0.45,6}{\draw[gray!55,thin] (\x,0) -- (\x,-3.82);}
\end{tikzpicture}
\caption{The two-surface construction behind the lensmaker's equation. The first surface alone would image $S$ at $P'$; here the rays inside the glass are still diverging, so $P'$ is virtual and lies to the \emph{left}, giving $s_{i1} < 0$. Whether real or virtual, $P'$ is the object seen by the second surface, at distance $s_{o2} = d - s_{i1}$ from $V_2$.}
\label{fig:lensmaker}
\end{figure}

\emph{First surface.} Light crosses from $n_m$ into $n_l$ over radius $R_1$, so the paraxial single-surface equation reads
\begin{equation}
    \frac{n_m}{s_{o1}} + \frac{n_l}{s_{i1}} = \frac{n_l - n_m}{R_1}.
    \label{eq:lens-surf1}
\end{equation}

\emph{Handing off to the second surface.} The image $P'$ that the first surface would form is the object the second surface sees --- whether or not the light ever physically converges there. Measuring from $V_2$, a distance $d$ farther along the axis,
\begin{equation}
    s_{o2} = d - s_{i1}.
\end{equation}
This one relation carries the whole sign convention. In Figure~\ref{fig:lensmaker} the rays inside the glass still diverge, so $P'$ is virtual, $s_{i1} < 0$, and $s_{o2} = d + |s_{i1}| > 0$: a perfectly ordinary real object for the second surface. Had the first surface instead converged the light to a real image lying beyond $V_2$, then $s_{i1} > d$ would give $s_{o2} < 0$, correctly flagging a virtual object. No case analysis is needed --- the formula handles both.

\emph{Second surface.} Light crosses back from $n_l$ into $n_m$ over radius $R_2$:
\begin{equation}
    \frac{n_l}{d - s_{i1}} + \frac{n_m}{s_{i2}} = \frac{n_m - n_l}{R_2}.
    \label{eq:lens-surf2}
\end{equation}

\emph{Adding the two.} Summing \eqref{eq:lens-surf1} and \eqref{eq:lens-surf2}, the right-hand sides combine into a single factor of $(n_l - n_m)$, while the two terms carrying $s_{i1}$ --- the only trace of the intermediate image --- collapse over a common denominator:
\begin{equation}
    \frac{n_l}{s_{i1}} + \frac{n_l}{d - s_{i1}}
    = n_l\,\frac{(d - s_{i1}) + s_{i1}}{s_{i1}(d - s_{i1})}
    = \frac{n_l\, d}{s_{i1}(d - s_{i1})}.
\end{equation}
The intermediate image distance therefore survives only in a single term, and only multiplied by the thickness:
\begin{equation}
    \frac{n_m}{s_{o1}} + \frac{n_m}{s_{i2}} + \frac{n_l\, d}{s_{i1}(d - s_{i1})}
    = (n_l - n_m)\left(\frac{1}{R_1} - \frac{1}{R_2}\right).
    \label{eq:thick-lens}
\end{equation}
Nothing has been approximated beyond the paraxial assumption already built into each surface equation: \eqref{eq:thick-lens} is the exact thick-lens result.

\emph{The thin-lens limit.} Letting $d \to 0$ kills the middle term outright, and with it every reference to $s_{i1}$ --- which is why a thin lens can be characterized without ever locating the intermediate image. Writing $s_o \equiv s_{o1}$ and $s_i \equiv s_{i2}$, both now measured from the single plane the lens has collapsed onto,
\begin{equation}
    \frac{1}{s_o} + \frac{1}{s_i} = \frac{n_l - n_m}{n_m}\left(\frac{1}{R_1} - \frac{1}{R_2}\right).
\end{equation}
The right-hand side involves only the lens and its surroundings, not the object, so it is a constant of the lens; setting $s_o \to \infty$ makes $s_i$ equal to that constant's reciprocal, which is precisely the definition of the focal length. Specializing to a lens of index $n$ in air ($n_m = 1$, $n_l = n$):
\begin{formula}[Lensmaker's equation]
\begin{equation}
    \frac{1}{f} = (n-1)\left(\frac{1}{R_1} - \frac{1}{R_2}\right).
\end{equation}
\end{formula}
\begin{formula}[Thin-lens (Gaussian) equation]
\begin{equation}
    \frac{1}{s_o} + \frac{1}{s_i} = \frac{1}{f}.
\end{equation}
\end{formula}
Here $f_o = f_i \equiv f$, since the medium is now common to both sides. Measuring distances instead from the focal points, with $x_o = s_o - f$ and $x_i = s_i - f$, the same content takes the compact \emph{Newtonian form} $x_o x_i = f^2$.

The reciprocal $P = 1/f$, in diopters when $f$ is in metres, is the lens \textbf{power}; thin lenses in contact simply add, $1/f = 1/f_1 + 1/f_2$, which is why prescriptions are written in diopters rather than focal lengths.

% ============================================================
\section{Mirrors}
% ============================================================

A spherical mirror can be treated exactly with no more than the law of reflection. A ray parallel to the axis striking at angle $\theta$ (measured at the centre of curvature) crosses the axis at
\begin{equation}
    f(\theta) = R - \frac{R}{2\cos\theta}
    \;\xrightarrow[\ \theta \to 0\ ]{}\; \frac{R}{2},
\end{equation}
so paraxially $f = R/2$, and the thin-lens equation $1/s_o + 1/s_i = 1/f$ applies unchanged. The residual $\theta$-dependence is spherical aberration in its purest form: as Figure~\ref{fig:mirror-aberr} shows, marginal rays cross the axis nearer the mirror than paraxial ones, so a sphere has no single focus.

\begin{figure}[H]
\centering
\begin{tikzpicture}[>=Stealth, scale=0.95]
  \def\R{4}
  % mirror arc (concave toward the left, centre of curvature at the origin)
  \draw[very thick] plot[domain=-48:48,samples=90] ({\R*cos(\x)},{\R*sin(\x)});
  \draw[dashed] (-1.0,0) -- (6.4,0);
  % incoming parallel rays, each reflected to its own axis crossing at R/(2 cos theta)
  \foreach \h/\col in {0.6/blue!70!black, 1.7/green!45!black, 2.7/red!75!black}{
    \pgfmathsetmacro{\th}{asin(\h/\R)}
    \pgfmathsetmacro{\xh}{\R*cos(\th)}
    \pgfmathsetmacro{\xc}{\R/(2*cos(\th))}
    \draw[thick,\col,->] (6.3,\h) -- ({\xh+0.05},\h);
    \draw[thick,\col,->] (\xh,\h) -- (\xc,0);
    \filldraw[\col] (\xc,0) circle (1.5pt);
  }
  \filldraw (0,0) circle (1.8pt) node[below left] {\small$C$};
  \filldraw (\R,0) circle (1.8pt) node[below right] {\small$V$};
  \node[below] at (2.0,-0.12) {\small$F = R/2$};
  % brace marking the spread of crossing points
  \draw[thick,gray,|<->|] (2.0,-0.95) -- (2.72,-0.95);
  \draw[gray,thin] (2.36,-1.0) -- (1.55,-1.5);
  \node[gray,align=center,left] at (1.5,-1.55) {\small spread of foci:\\ \small spherical aberration};
\end{tikzpicture}
\caption{A spherical mirror has no single focus. Rays striking farther off axis cross the axis at $f(\theta) = R - R/(2\cos\theta)$, which drifts away from the paraxial value $R/2$ as the ray height grows.}
\label{fig:mirror-aberr}
\end{figure}

The cure is to abandon the sphere. Each conic section images exactly one conjugate pair stigmatically: a \emph{paraboloid} for an object at infinity, since by the focus--directrix property every axis-parallel ray has the same OPL to the focus --- hence parabolic telescope primaries; an \emph{ellipsoid} between its two finite foci; a \emph{hyperboloid} for one real and one virtual focus. Spheres are exact for no pair at all, and survive only because they are enormously easier to grind.

% ============================================================
\section{Magnification}
% ============================================================

Similar triangles on the ray through the lens centre --- undeviated, because a thin lens there is a parallel-faced slab of vanishing thickness --- give the transverse magnification:
\begin{formula}[Transverse magnification]
\begin{equation}
    m = \frac{h_i}{h_o} = -\frac{s_i}{s_o}.
\end{equation}
\end{formula}
Positive $m$ is upright, negative inverted, and for a compound system the magnifications simply multiply, $m_{\text{tot}} = \prod_k m_k$. Note that a real object imaged really ($s_o, s_i > 0$) forces $m < 0$: such an image is \emph{always} inverted, by geometry rather than by convention.

% ============================================================
\section{Compound Systems}
% ============================================================

Compound systems need no new machinery: apply the single-element equation in sequence, letting the image of element $k$ be the object of element $k+1$. With a separation $L$ between them,
\begin{equation}
    s_{o,k+1} = L - s_{i,k}.
\end{equation}
A negative result here is a \textbf{virtual object} --- converging rays intercepted by the next element before they ever formed an image --- and the sign convention handles it without modification.

% ============================================================
\section{Summary}
% ============================================================

\begin{fact}[Paraxial toolkit]
\begin{align}
    \text{Snell:}\quad n_1\sin\theta_1 &= n_2\sin\theta_2,
    &\text{Critical angle:}\quad \sin\theta_c &= \frac{n_2}{n_1}, \\
    \text{Single surface:}\quad \frac{n_1}{s_o} + \frac{n_2}{s_i} &= \frac{n_2-n_1}{R},
    &\text{Mirror:}\quad f &= \frac{R}{2}, \\
    \text{Gaussian:}\quad \frac{1}{s_o} + \frac{1}{s_i} &= \frac{1}{f},
    &\text{Newtonian:}\quad x_o x_i &= f^2, \\
    \text{Lensmaker:}\quad \frac{1}{f} &= (n-1)\left(\frac{1}{R_1}-\frac{1}{R_2}\right),
    &\text{Magnification:}\quad m &= -\frac{s_i}{s_o}.
\end{align}
\end{fact}

The unifying lesson mirrors the one that closed the diffraction chapter. There, every pattern was a Fourier transform of the aperture; here, every imaging equation is Fermat's principle evaluated on a particular surface, and the paraxial approximation is what makes the ray height cancel so that an image exists at all.
